\documentclass[11pt]{article}
\usepackage[margin=1.1in]{geometry}
\usepackage{amsmath,amssymb,amsthm}
\usepackage[colorlinks=true,linkcolor=blue,citecolor=blue]{hyperref}

\newtheorem{theorem}{Theorem}[section]
\newtheorem{lemma}[theorem]{Lemma}
\newtheorem{proposition}[theorem]{Proposition}
\newtheorem{corollary}[theorem]{Corollary}
\newtheorem{definition}[theorem]{Definition}
\newtheorem{remark}[theorem]{Remark}
\newtheorem{problem}[theorem]{Problem}

\newcommand{\R}{\mathbb{R}}
\newcommand{\E}{\mathbb{E}}
\newcommand{\PP}{\mathbb{P}}
\newcommand{\emp}{\bar\mu_N}
\newcommand{\vol}{\operatorname{vol}}
\newcommand{\dist}{\operatorname{dist}}
\newcommand{\Hess}{\operatorname{Hess}}
\newcommand{\diam}{\operatorname{diam}}
\newcommand{\Acal}{\mathcal A}
\newcommand{\Zcal}{\mathcal Z}
\newcommand{\Ncal}{\mathcal N}

\title{The interpolation law of the uniform doubling condition\\
for curved truncations}
\author{Yitzchak Shmalo\\
\small Einstein Institute of Mathematics, The Hebrew University of Jerusalem, Jerusalem, Israel\\
\small }
\date{June 2026}

\begin{document}
\maketitle

\begin{abstract}
In companion papers the critical scaling of the uniform doubling condition of
[Shmalo, \emph{Ann.\ Math.\ Artif.\ Intell.} (2024)] for i.i.d.\ samples was
determined to be $\Theta\bigl((\log N/N)^{2/(2n-j)}\bigr)$ at every transversal
boundary stratum of dimension $j$, first in the orthogonal model and then for
all simplicial (half-space) truncations in general position, the upper bound
resting on a witness-box section lemma that diagonalizes the active
constraints.  Those arguments use the flatness of the truncating surfaces in
two places: the level constraint of a half-space is linear, and the slice of a
half-space by a coordinate hyperplane is again a half-space.  Here we remove
both uses for a \emph{fixed} smooth strictly convex truncation family.  Let
$G=K\cap\bigcap_{i=1}^k\{g_i\le0\}$ with $K$ and the $g_i$ fixed and $C^2$, the
surfaces $\{g_i=0\}$ and $\partial K$ strictly convex and meeting transversally
in general position along a $j$-dimensional stratum.  We prove that the cap of
depth $h$ at the stratum has volume $\asymp h^{(2n-j)/2}$ with the same exponent
as in the flat case, the constant now depending on the joint second fundamental
form of $K$ and the $g_i$ but bounded above and below for fixed $C^2$ data; and
that the slice of $G$ by a level set of one constraint is an admissible
$(n-1,k-1)$ curved configuration, with curvature bounds depending on the fixed
configuration and uniform along the sliced stratum, closing the
slicing induction.  Consequently the interpolation law
$\Theta\bigl((\log N/N)^{2/(2n-j)}\bigr)$ holds for the fixed curved family.
The mechanism that survives curvature is the steep/shallow trichotomy at the
angular threshold $\sqrt h$: a strictly convex constraint contributes a factor
$h$ through its linear part when the slab direction is transversal to its
normal, and degrades continuously to a factor $\sqrt h$ exactly when the slab
direction approaches that normal, at which point a lower-dimensional stratum
becomes active and the exponent law reports the correct larger value.  We
isolate honestly what is and is not obtained: for a fixed family the constants
are uniform over the stratum and over directions in any fixed compact subcone
of the active normal cone, which suffices for the fixed-family interpolation
law; a curvature-band-uniform lower/slab constant near the multiply-tangent
directions remains open, and we state the residual estimate precisely.  We also
prove a conditional uniform theorem: if iterated slices keep a uniform
curvature floor and all simultaneously projected near-tangent normals have a
uniform Gram lower bound, then the slab constants and the interpolation law are
uniform over that resolved subfamily.  As an immediate consequence, every
finite atlas of fixed curved configurations, and every compact atlas whose
resolved constants vary continuously and stay positive, has a uniform law with
atlas-dependent constants.  Resolved geometry is also stable under
\(C^2\)-small perturbations, so any compact \(C^2\)-continuous resolved atlas is
uniform by a finite subcover.  The same statement is locally uniform on every
compact stage of a \(\sigma\)-compact resolved atlas, and becomes globally
uniform exactly when the resolved constants have a positive global floor.  For
one curved truncation, the resolved hypotheses are automatic under quantitative
splaying, and also under an explicit curvature-dominance inequality that allows
controlled oblique intersections; in particular, nonobtuse one-cut contact is
uniform.
\end{abstract}

\section{Setting and results}\label{sec:intro}

We keep the framework of the companion papers
\cite{companion2,companion3,companion4,companion11,companion6}.  Let $K\subset\R^n$
be a convex body, $\mu$ a probability measure with density bounded above and
below on $K$, and $T=\{X_1,\dots,X_N\}$ i.i.d.\ samples with empirical measure
$\emp$.  For a unit vector $u$, center $s$, width $w$ and a truncation region
$G$, the truncated slab is
$S_w=\{x:|\langle x-s,u\rangle|\le\frac w2\}\cap G$, and the uniform doubling
condition (UDC) with constants $(\kappa,\delta,\sigma,\beta,\ell)$ requires
$\emp(S_{\ell\kappa^{i+1}})\ge\min\{\delta,(1+\sigma)\emp(S_{\ell\kappa^{i}})\}$
for every truncated slab generated by the family and every grid scale
$\ell\kappa^i<\beta$.  By \cite{companion3,companion4,companion11,companion6} the
critical scaling of $\ell$ for half-space truncation families in general
position is $\Theta\bigl((\log N/N)^{2/(2n-j)}\bigr)$, where $j$ is the
dimension of the lowest transversal boundary stratum; the lower bound comes
from a second-moment count over $\asymp(N/\log N)^{j/(2n-j)}$ congruent caps
along the stratum, and the upper bound from a deterministic slab-volume
estimate $\vol(G\cap\{|\langle y-x,u\rangle|\le L\})\gtrsim L^{(2n-j)/2}$ fed
into the proof skeleton of \cite{companion2}.

The immediately preceding companion \cite{companion9} removes the orthogonality
and roundness of the orthogonal model, treating a $C^2$ strictly convex body $K$
cut by half-spaces in general position, and reduces the interpolation law to the
flat-truncation case; it leaves open exactly the configurations whose truncating
\emph{surfaces} are themselves curved.  This note treats that case.

\begin{definition}[Admissible curved configuration]\label{def:config}
Fix $n\ge2$ and $0\le k\le n-2$, and put $j:=n-1-k$.  Let
$K\subset\R^n$ be a compact convex body with $C^2$ boundary whose principal
curvatures lie in $[c_K,C_K]$, $0<c_K\le C_K$, and let
$g_1,\dots,g_k\in C^2(\R^n)$ be convex, with
\begin{equation}\label{eq:reg}
0<\gamma_0\le|\nabla g_i|\le\Gamma_0,
\qquad
c_g\,I\preceq\Hess g_i\preceq C_g\,I
\quad\text{on a neighborhood of }G,
\end{equation}
($0<c_g\le C_g$), so that each $\{g_i=0\}$ is a $C^2$ strictly convex
hypersurface.  Write
$G:=K\cap\bigcap_{i=1}^k\{g_i\le0\}$ and let $n_i:=\nabla g_i/|\nabla g_i|$ be
the outward unit normal of the $i$-th surface.  We call the data
\emph{admissible} if $G$ has nonempty interior, contains a ball of a fixed
radius $r_0>0$, and the surfaces meet \emph{transversally in general position
along a $j$-dimensional stratum}
\[
E:=\partial K\cap\bigcap_{i=1}^k\{g_i=0\},
\]
a compact $C^2$ $j$-manifold, meaning: at every $p\in E$ the $k+1$ outward unit
normals $\nu_K(p),n_1(p),\dots,n_k(p)$ have Gram matrix
$M(p)$ with smallest eigenvalue $\ge\gamma>0$, uniformly in $p\in E$.  We call
$(c_K,C_K,c_g,C_g,\gamma_0,\Gamma_0,\gamma,r_0)$ the \emph{structure constants}.
\end{definition}

The simplicial model of \cite{companion11} is the degenerate case $C_g=0$,
$\Hess g_i\equiv0$ (linear $g_i$), with $K=B^n$.  Our standing example is a
ball truncated by two spheres, $K=B^n$ and
$g_i(x)=|x-z_i|^2-\rho_i^2$ with the $z_i$ placed so that the two spheres and
$\partial K$ meet transversally along an $(n-3)$-sphere.

For $p\in E$ the maximum of a linear functional over $G$ is governed, to first
order, by the active constraints at $p$.  As in \cite{companion11} the relevant
extremal directions are those of the \emph{deepest corner}: $u$ for which the
maximizer of $\langle\cdot,u\rangle$ over $G$ is the single point $p\in E$ with
\emph{all $k+1$ Karush--Kuhn--Tucker multipliers strictly positive}.  These
form the relative interior of the normal cone of $E$ at $p$,
\[
\mathcal U_p:=\operatorname{relint}\Bigl\{\lambda_K\nu_K(p)+\sum_{i=1}^k\lambda_i n_i(p):
\ \lambda_K,\lambda_i\ge0\Bigr\}\cap S^{n-1}.
\]

\begin{theorem}[Cap volume at a curved stratum]\label{thm:cap}
Let the data be admissible with structure constants as above, and fix
$\theta\in(0,1]$.  There are $h_1>0$ and $b_2>0$ depending on the fixed
configuration data, and $b_1=b_1(\theta)>0$ depending in addition on $\theta$,
such that for every $p\in E$, every $u\in\mathcal U_p$ at sphere-distance at
least $\theta$ from $\partial\mathcal U_p$, and every $0<h\le h_1\theta^2$,
\[
b_1(\theta)\,h^{\frac{2n-j}2}\;\le\;
\mu\bigl(\{x\in G:\ \langle x,u\rangle\ge\tau_{\max}(u)-h\}\bigr)
\;\le\;b_2\,h^{\frac{2n-j}2},
\qquad\tau_{\max}(u):=\max_{x\in G}\langle x,u\rangle .
\]
The upper constant $b_2$ is uniform over the whole cone (no $\theta$); the lower
constant degrades as $\theta\to0$ ($b_1(\theta)\gtrsim\theta^{\,k+j}$), staying
positive for each fixed $\theta$.  In particular, for $u$ ranging over any
compact subcone of $\mathcal U_p$ (fixed $\theta$), the cap exponent is exactly
$(2n-j)/2$, independent of the curvatures, with constants determined by the
fixed configuration data and $\theta$.
\end{theorem}

\begin{theorem}[Section lemma and slab bound]\label{thm:slab}
For each fixed configuration satisfying Definition~\ref{def:config} there are $a_1>0$ and
$L_0>0$, depending on that configuration, such that for
every $x\in G$, every unit $u$, and every $0<L\le L_0$,
\[
\vol\Bigl(G\cap\{y:\ |\langle y-x,u\rangle|\le L\}\Bigr)\;\ge\;a_1\,L^{\frac{2n-j}2}.
\]
\end{theorem}

\begin{corollary}[Interpolation law, fixed curved family]\label{cor:law}
Let the data be admissible.  Then the critical scaling of the uniform doubling
condition for the (single fixed) curved truncation family
$\{g_1,\dots,g_k\}$ on $K$ is
\[
\ell_N\;\asymp\;\Bigl(\frac{\log N}{N}\Bigr)^{2/(2n-j)},
\qquad j=n-1-k,
\]
both bounds holding with probability $1-o(1)$, exponent and logarithm exact.
The implied constants depend on the fixed configuration data (equivalently, on
the structure constants together with the nondegeneracy constants of this
family), but not on $N$.
\end{corollary}

Section~\ref{sec:cap} proves Theorem~\ref{thm:cap}; Section~\ref{sec:slice}
proves the slicing induction behind Theorem~\ref{thm:slab} and assembles
Corollary~\ref{cor:law}; Section~\ref{sec:uniform} states precisely what
uniformity is and is not obtained, isolates the residual open estimate, and
proves a conditional uniform theorem under resolved multi-tangent geometry.

\section{The cap-volume proposition with curved constraints}\label{sec:cap}

Throughout this section fix $p\in E$ and $u\in\mathcal U_p$ at sphere-distance
$\ge\theta$ from $\partial\mathcal U_p$.  After a fixed rotation and translation
we place $p$ at the origin with $u=e_1$ pointing in the direction of increasing
$\tau=\langle\cdot,u\rangle$; thus $\tau_{\max}(u)=0$ and the cap of depth $h$ is
\[
C(h)=\{x\in G:\ x_1\ge-h\}.
\]
Write $x=(x_1,\xi)$ with $\xi\in\R^{n-1}$ the slab coordinates (the orthogonal
complement of $u$).

\subsection{The local floor}\label{ss:floor}

At $p=0$ the $k+1$ active surfaces are $\partial K$ and $\{g_i=0\}$,
$i=1,\dots,k$.  Each is a $C^2$ graph over its tangent plane at $0$; since all
$k+1$ outward normals lie in the open half-space $\{x_1>0\}$-side (because
$u=e_1$ is a strictly positive combination of them), each surface is, near $0$,
a graph $x_1=\psi_\bullet(\xi)$ over the slab, with
\begin{equation}\label{eq:graphs}
\psi_K(\xi)=-\langle a_K,\xi\rangle-\tfrac12\xi^\top Q_K\,\xi+O(|\xi|^3),
\qquad
\psi_i(\xi)=-\langle a_i,\xi\rangle-\tfrac12\xi^\top Q_i\,\xi+O(|\xi|^3),
\end{equation}
where $a_\bullet\in\R^{n-1}$ is the slab-gradient of the surface at $0$ and
$Q_\bullet\succeq0$ is its second fundamental form in the slab coordinates
(positive semidefinite by convexity of $K$ and of the $g_i$; positive definite
on the $j$ stratum directions for $K$, see below).  The $O(|\xi|^3)$ remainders
are uniform on $E$ for fixed $C^2$ data.  The region $G$ near $0$ is
\begin{equation}\label{eq:feasible}
C(h)=\Bigl\{(x_1,\xi):\ -h\le x_1\le\min\bigl(\psi_K(\xi),\,\psi_1(\xi),\dots,\psi_k(\xi)\bigr)\Bigr\}.
\end{equation}
The \emph{floor} is the function $\Phi(\xi):=\min_\bullet\psi_\bullet(\xi)$, and
the cap is the region between the height $-h$ and the floor wherever the floor
exceeds $-h$.

The linear data $a_\bullet$ encode the angle between $u$ and the surface
normals.  Indeed if $\nu$ is the outward unit normal of a surface
$\{x_1=\psi(\xi)\}$ at $0$, then writing $\nu=(\nu_1,\nu')$ with
$\nu_1=\langle\nu,u\rangle$ one has $a=-\nu'/\nu_1$, hence
\begin{equation}\label{eq:beta}
|a_\bullet|=\frac{|\nu'_\bullet|}{\langle\nu_\bullet,u\rangle}
=\frac{\sin\angle(\nu_\bullet,u)}{\cos\angle(\nu_\bullet,u)}=\tan\angle(\nu_\bullet,u),
\qquad
\beta_\bullet:=\frac1{\sqrt{1+|a_\bullet|^2}}=\langle\nu_\bullet,u\rangle .
\end{equation}
The quantity $\beta_\bullet=\langle\nu_\bullet,u\rangle\in(0,1]$ is the cosine of
the angle between $u$ and the $\bullet$-normal, and
$m_\bullet:=|a_\bullet|=\tan\angle(\nu_\bullet,u)=\sqrt{1-\beta_\bullet^2}/\beta_\bullet$
is the slope of the surface graph in the slab coordinates.  The slope governs
the pinning: along the steepest slab direction
$\hat a_\bullet:=a_\bullet/|a_\bullet|$ the floor $\psi_\bullet$ decreases at
unit rate $m_\bullet$, so the constraint $x_1\le\psi_\bullet$ at height
$x_1\ge-h$ confines that coordinate to a window of length
$\asymp h/m_\bullet$.  A \emph{steep} surface (slope $m_\bullet$ bounded, $u$
transversal to the normal) has a window of length $\asymp h$; a \emph{shallow}
one ($m_\bullet$ large, $u$ near the surface's tangent plane,
$\beta_\bullet\to0$) has a shorter window, but the surface curvature then caps
the relevant extent at $\sqrt h$.  This is the curved analogue of the
$u_i$-window $h/(8ku_i)$ of \cite{companion11}, where $u_i=\langle u,e_i\rangle$
is exactly the cosine $\beta_i$ for the orthogonal model.  We record the two
structural facts.

\begin{lemma}[Curvature is positive definite on the stratum]\label{lem:pd}
Let $V\subseteq\R^{n-1}$ be the slab-image of $T_pE$, the tangent space of the
stratum, $\dim V=j$.  Then there is a structure constant $c'_K\asymp c_K>0$ such
that for every $p\in E$ and every admissible $u$,
\[
Q_K\big|_V\;\succeq\;c'_K\,I_V,
\]
so the floor is strictly convex along the stratum directions.  Consequently the
joint tangential form $Q_K\big|_V+\sum_i Q_i\big|_V\succeq c'_K I_V$ as well,
since each $Q_i\succeq0$ by convexity of $g_i$.
\end{lemma}

\begin{proof}
$\partial K$ has second fundamental form $\mathrm{I\!I}_K\succeq c_K I$ on its
tangent plane $T_p\partial K$.  Writing $\partial K$ near $p$ as a graph
$x_1=\psi_K(\xi)$ over the slab $u^\perp$, tilted from $T_p\partial K$ by the
angle $\alpha=\angle(\nu_K,u)$, the graph Hessian is
$Q_K=\mathrm{I\!I}_K/\cos^3\!\alpha$ in the tilt direction and
$\mathrm{I\!I}_K/\cos\alpha$ transversally, hence $Q_K\succeq\mathrm{I\!I}_K
\succeq c_K I$ on $u^\perp$ (tilting only increases the graph curvature).
Restricting to $V$, the slab projection $T_pE\to V$ is a linear isomorphism with
singular values bounded above by $1$, so $Q_K|_V\succeq c'_K I_V$ with a
structure constant $c'_K\asymp c_K$ that does \emph{not} degrade as $u$
approaches a facet (the tilt makes $Q_K$ larger, not smaller).  The last
sentence is immediate from $Q_i\succeq0$.
\end{proof}

\begin{lemma}[Transversal lower bound on the corner angle]\label{lem:angle}
There is a structure constant $\beta_*>0$ (depending on $\gamma$ and on $n,k$)
such that for every $p\in E$ and every $u\in\mathcal U_p$ at sphere-distance
$\ge\theta$ from $\partial\mathcal U_p$,
\[
\beta_\bullet=\langle\nu_\bullet,u\rangle\;\ge\;\beta_*\,\theta
\qquad\text{for each active normal }\nu_\bullet\in\{\nu_K,n_1,\dots,n_k\}.
\]
\end{lemma}

\begin{proof}
Write $u=\sum_\bullet\lambda_\bullet\nu_\bullet$ with $\lambda_\bullet\ge0$ and,
since $|u|=1$, $\lambda^\top M\lambda=1$, where $M\succeq\gamma I$ is the Gram
matrix.  Hence $|\lambda|\le\gamma^{-1/2}$, and
$\langle\nu_\bullet,u\rangle=(M\lambda)_\bullet$.  If $u$ lies at
sphere-distance $\ge\theta$ from $\partial\mathcal U_p$, then $u$ is not within
$\theta$ of any facet $\{\lambda_\bullet=0\}$ of the cone; equivalently each
$\lambda_\bullet$ is bounded below by a multiple of $\theta$ (the facet
distance is comparable to the smallest multiplier, with comparison constants
fixed by $M\in[\gamma,\Gamma_0^2]$).  Thus
$(M\lambda)_\bullet\ge\gamma\lambda_\bullet\ge\beta_*\theta$.  Consequently the
slope is bounded, $m_\bullet=\sqrt{1-\beta_\bullet^2}/\beta_\bullet\le
1/\beta_\bullet\le1/(\beta_*\theta)$.
\end{proof}

\subsection{The witness box}\label{ss:box}

We now run the witness-box construction of \cite{companion11} with the linear
constraint $x_i\le0$ replaced by the curved floor \eqref{eq:graphs}, and the
orthogonality $\langle w,\tilde u\rangle=0$ replaced by working directly in the
slab coordinates $\xi$, in which the height $x_1$ is the free fibre direction.
Split the active surfaces by the slope threshold at the angular scale
$\sqrt h$: with $C_*:=C_g+C_K$ bounding all floor curvatures,
\begin{equation}\label{eq:trichotomy}
\Acal:=\bigl\{\bullet:\ m_\bullet\le C_*^{-1/2}h^{-1/2}\bigr\},
\qquad
\Zcal:=\bigl\{\bullet:\ m_\bullet> C_*^{-1/2}h^{-1/2}\bigr\}.
\end{equation}
We call $\Acal$ the \emph{steep} surfaces (bounded slope: $u$ transversal to the
normal, $\beta_\bullet$ not too small) and $\Zcal$ the \emph{shallow}
(slope $\gtrsim h^{-1/2}$: $u$ near the surface's tangent plane,
$\beta_\bullet\lesssim\sqrt h$).  By Lemma~\ref{lem:angle},
$m_\bullet\le1/(\beta_*\theta)$ for $u$ at facet-distance $\ge\theta$, so every
active surface is steep, $\Zcal=\varnothing$, once
$1/(\beta_*\theta)\le C_*^{-1/2}h^{-1/2}$, i.e.\ once
$h\le\beta_*^2\theta^2/C_*=:h_1\theta^2$.  This is the depth range asserted in
Theorem~\ref{thm:cap}; we carry out the argument there.  The role of the
trichotomy across the transition $\Zcal\ne\varnothing$ is deferred to
Section~\ref{sec:uniform}.

Fix $h\le h_1\theta^2$, so every $m_\bullet\le C_*^{-1/2}h^{-1/2}$, equivalently
the pinned window $h/m_\bullet\ge\sqrt{h/C_*}\cdot$(a constant) and the linear
part of each $\psi_\bullet$ dominates its quadratic part at scale $h$.  Let $V$
be the slab-image of $T_pE$ as in Lemma~\ref{lem:pd}, and let
$W:=V^\perp\subseteq\R^{n-1}$, $\dim W=k$.  The slab-projected constraint normals
$P_u n_1,\dots,P_u n_k$ span $W$; by transversality their Gram is bounded below
by a structure constant, so there is a \emph{biorthogonal (dual) frame}
$f_1,\dots,f_k\in W$,
\begin{equation}\label{eq:dual}
\langle f_i,\,P_u n_{i'}\rangle=\delta_{ii'},
\qquad |f_i|\le F_0\ \text{(a structure constant)},
\end{equation}
bounded because the Gram is bounded away from singular.  This frame plays the
role of the orthogonality $\langle w,\tilde u\rangle=0$ of \cite{companion11}: it
provides coordinates $s_i:=\langle\xi,P_u n_i\rangle$ for $\xi\in\R^{n-1}$ in
which the $i$-th constraint depends, to leading order, only on $s_i$, namely
$\psi_i(\xi)=-s_i-\frac12\xi^\top Q_i\xi+O(|\xi|^3)$ (the linear part of
$\psi_i$ is $-\langle a_i,\xi\rangle$ with $a_i=-P_u n_i/\beta_i$, so
$-\langle a_i,\xi\rangle=s_i/\beta_i$; we absorb the bounded factor $1/\beta_i$
into the unit by rescaling $s_i$, keeping $|s_i|\asymp\dist$ to the $i$-th
surface).  Write $\xi=(y,z)$ with $y\in V$ the stratum coordinates and
$z=\sum_i s_i f_i\in W$ the pinned coordinates.

For the cap-volume bounds of Theorem~\ref{thm:cap} we work at fixed
facet-distance $\theta$, where by Lemma~\ref{lem:angle} every slope is bounded
by the \emph{fixed} constant $M_0:=1/(\beta_*\theta)$ and every cosine
$\beta_\bullet\ge\beta_*\theta$.  Pin each constraint at scale $h$: define the
\emph{curved witness box} $\widehat W$ in coordinates $(x_1,y,z)$ by
\begin{equation}\label{eq:cwbox}
|y|\le\eta\sqrt h,
\qquad
-\eta\,h\le s_i\le 0\ \ (i=1,\dots,k),
\qquad
-h\le x_1\le-\tfrac h2,
\end{equation}
for a small structure constant $\eta>0$ (we will need $\eta\le\beta_*\theta$ and
$\eta$ small in absolute constants); $s_i\le0$ is the inward side of the $i$-th
constraint.  The pinned windows have length $\eta h\asymp h$ --- this is the
range that produces the factor $h$ per constraint --- so $|z|\le F_0\sum_i|s_i|
\le kF_0\eta h\le c\,h$, while the stratum coordinates carry $|y|\le\eta\sqrt h$.
Thus $|z|=O(h)$ but $|y|=O(\sqrt h)$, and $|\xi|^2=|y|^2+|z|^2\le\eta^2 h+c^2h^2
\le c_1\eta^2 h$ for $h$ small.

\paragraph{$\widehat W$ lies in the cap.}
We verify $x_1\le\Phi(\xi)=\min_\bullet\psi_\bullet(\xi)$ on $\widehat W$.  The
curvature term of each $\psi_\bullet$ is
$\frac12|\xi^\top Q_\bullet\xi|\le\frac12 C_*c_1\eta^2 h\le\frac h8$ ($\eta$
small), and the cubic remainder is $O((\eta^2 h)^{3/2})\le\frac h8$ ($h$ small).
For the linear term $-\langle a_\bullet,\xi\rangle=-\langle a_\bullet|_V,y\rangle
-\langle a_\bullet|_W,z\rangle$:
\begin{itemize}
\item \emph{Constraint $\bullet=i$.}  By the dual frame \eqref{eq:dual},
$\langle a_i,\xi\rangle=-s_i/\beta_i$ has stratum part zero
($a_i|_V$ contributes through $y$ only at quadratic order, absorbed above, since
$P_un_i\perp V$) and pinned part $-s_i/\beta_i\ge0$ for $s_i\le0$.  Hence
$-\langle a_i,\xi\rangle=s_i/\beta_i\le0$: the pinning term only \emph{helps}
$x_1\le\psi_i$.  Its size is $|s_i|/\beta_i\le\eta h/(\beta_*\theta)\le h$ by
$\eta\le\beta_*\theta$, well within budget.
\item \emph{Body $\bullet=K$.}  Since the contact point is the deepest corner,
$0$ maximizes $\tau$ over $K$ along the whole stratum, so $a_K|_V=0$.  The pinned
part has slope $\le M_0$, so $|\langle a_K|_W,z\rangle|\le M_0|z|\le M_0c\,h
\le\frac h8$ ($\eta$ small, since $|z|\le c\eta h$ carries the smallness).
\end{itemize}
Collecting, $\psi_\bullet(\xi)\ge-\frac h8-\frac h8-\frac h8\ge-\frac h2\ge x_1$
on $\widehat W$, so $\widehat W\subseteq C(h)$.  The essential point is that the
pinned coordinates sit at scale $h\ll\sqrt h$, so they contribute nothing at the
curvature scale; only the stratum coordinates $y$ reach scale $\sqrt h$, and
there the body is critical ($a_K|_V=0$) and contributes only its quadratic drop,
itself $\le\frac h8$.

\paragraph{The pinning is genuinely linear.}
This is the heart of the curved estimate.  Along the $i$-th constraint
direction the floor is $\psi_i=s_i/\beta_i-\frac12\xi^\top Q_i\xi+O(|\xi|^3)$.
The \emph{linear} coefficient of $s_i$ is $1/\beta_i\ge1$, bounded below, so the
feasibility $\psi_i\ge-h$ permits $|s_i|\le\beta_i h+O(h^{3/2})\asymp h$ (the
window length used in \eqref{eq:cwbox}); the \emph{quadratic} term
$\frac12 Q_i s_i^2=O(C_g h^2)\ll h$ cannot change this.  This is the exact
analogue of the linear constraint $x_i\le0$ of \cite{companion11} pinning
$w_i\in[-h/(8ku_i),0]$, with $u_i=\beta_i$; curvature perturbs the window only by
a factor $1+O(h)$.  The mechanism that could change the exponent --- the slope
$m_i$ growing to scale $h^{-1/2}$, so the linear coefficient $1/\beta_i$ itself
reaches $h^{-1/2}$ and the window contracts toward $\sqrt h$ --- does not occur
at fixed $\theta$, where $\beta_i\ge\beta_*\theta$ is bounded below; it is the
subject of Section~\ref{sec:uniform}.

\paragraph{Volume.}
The $k$ pinned coordinates $s_i$ range over intervals of length $\eta h$; the
dual frame $\{f_i\}$ has bounded Gram, so the $z$-Jacobian is a structure
constant; the $j$ stratum coordinates $y$ range over a ball of radius
$\eta\sqrt h$; and the height fibre has length $\frac h2$.  Hence
\begin{equation}\label{eq:lower}
\mu(C(h))\;\ge\;\frac{\vol(\widehat W)}{\vol(B^n)}
\;\ge\;c_2\cdot \tfrac h2\cdot\prod_{i=1}^k(\eta h)\cdot(\eta\sqrt h)^{j}
\;=\;c_3\,h^{\,1+k+j/2}\;=\;c_3\,h^{\frac{2n-j}2},
\end{equation}
since $1+k+j/2=1+k+(n-1-k)/2=(n+1+k)/2=(2n-j)/2$.  Here $c_3\asymp\eta^{\,k+j}$
times structure constants; with the admissible choice $\eta=\beta_*\theta$ this
gives $b_1(\theta)\asymp\theta^{\,k+j}$, the stated $\theta$-dependence of the
lower constant (positive for each fixed $\theta>0$, degrading as
$\theta\to0$ --- consistent with the cap growing as $u$ nears
$\partial\mathcal U_p$).  This proves the lower bound of Theorem~\ref{thm:cap}.

\subsection{The upper bound}\label{ss:upper}

For the upper bound, by \eqref{eq:feasible} the cap is contained in
$\{-h\le x_1\le\Phi(\xi)\}\cap\{\Phi(\xi)\ge-h\}$.  Fix any single steep
constraint, say $K$ (steep for all admissible $h$).  Its graph gives
$x_1\le\psi_K(\xi)=-\langle a_K,\xi\rangle-\frac12\xi^\top Q_K\xi+O(|\xi|^3)$,
and the height fibre has length at most $\psi_K(\xi)+h$, so
\[
\mu(C(h))\;\le\;\frac1{\vol(B^n)}\int_{\{\Phi\ge-h\}}\bigl(\Phi(\xi)+h\bigr)_+\,d\xi
\;\le\;\frac1{\vol(B^n)}\int_{\{\psi_K\ge-h\}}\bigl(\psi_K(\xi)+h\bigr)\,d\xi .
\]
We bound the integration domain $\{\psi_K\ge-h\}\cap\bigcap_i\{\psi_i\ge-h\}$.
The body constraint $\psi_K\ge-h$, by strict convexity $Q_K|_V\succeq c'_K I_V$
(Lemma~\ref{lem:pd}) and $a_K|_V=0$, forces $\frac12 c'_K|y|^2\le h+o(h)$, i.e.
$|y|\le C\sqrt{h}$ in the stratum directions.  Each constraint $\psi_i\ge-h$,
whose linear coefficient $1/\beta_i\ge1$ in the pinned coordinate $s_i$ is
bounded below, forces $|s_i|\le\beta_i h+O(h^{3/2})\le 2h$ (using $\beta_i\le1$)
in the pinned directions.  The integrand $\psi_K(\xi)+h$ is $\le 2h$.  Hence,
with the bounded $z$-Jacobian of the dual frame,
\[
\mu(C(h))\;\le\;C\cdot\underbrace{(C\sqrt h)^{\,j}}_{y\text{-ball}}\cdot
\underbrace{\prod_{i=1}^k(2h)}_{s_i\text{-intervals}}\cdot\underbrace{(2h)}_{\text{fibre}}
\;=\;C'\,h^{j/2}\cdot h^{k}\cdot h\;=\;b_2\,h^{\frac{2n-j}2},
\]
the exponent being $j/2+k+1=(2n-j)/2$.  Here the pinned length $|s_i|\le2h$ is
curvature- and $\theta$-independent (it uses only $\beta_i\le1$ and the lower
bound $1/\beta_i\ge1$ on the linear coefficient), so $b_2$ depends only on the
structure constants --- in particular the upper bound is uniform over the whole
normal cone, not just at fixed facet-distance.  This proves
Theorem~\ref{thm:cap}.\qquad$\square$

\begin{remark}[Why the constant no longer factors]\label{rem:factor}
In the simplicial model the cap constant factored as a product of a
stratum-curvature factor (from $Q_K$) and one factor per truncation depending
only on its angle, because the constraints were mutually orthogonal and the
floor split as a sum of decoupled coordinates.  With curved $g_i$ the cross
blocks of the floor $\sum_\bullet\frac12\xi^\top Q_\bullet\xi$ couple the
stratum directions $y$ to the pinned directions $z$, and the dual frame
\eqref{eq:dual} (which inverts the joint Gram of all active normals) mixes the
truncation curvatures $Q_i$ with the stratum curvature $Q_K$.  Completing the
square --- shifting $y\mapsto y-Lz$ to remove the $y$--$z$ cross block, with
$L$ bounded because $Q_K|_V\succeq c'_K I_V$ is invertible --- moves $y$ by
$O(z)=O(h)\ll\sqrt h$ and changes the $z$-part at $O(h^2)\ll h$, so the product
volume count survives with a constant depending on the joint form
$Q_K|_V+\sum_iQ_i|_V$ rather than on $Q_K$ and the $Q_i$ separately.  The
exponent is untouched precisely because the shift is lower order at the cap
scale.
\end{remark}

\section{The slicing induction for curved slices}\label{sec:slice}

The cap-volume proposition of Section~\ref{sec:cap} handles the
\emph{stratum-direction} extremes ($u\in\mathcal U_p$).  As in
\cite{companion11}, the general direction is reduced to the stratum case by an
induction that slices off one constraint at a time.  In the simplicial model
the slice of a half-space by a coordinate hyperplane is again a half-space, so
the slice stays in class trivially.  Here the slice is a curved hypersurface
and we must verify it remains an admissible $(n-1,k-1)$ curved configuration
with curvatures bounded for the fixed configuration.

\subsection{The slice is in class}\label{ss:inclass}

Fix a constraint, say $g_{i_0}$, and a level $s\in\R$ near $0$.  Let
\[
\Sigma_s:=\{x:\ g_{i_0}(x)=s\},
\qquad
G^{(i_0)}_s:=G\cap\Sigma_s
=K\cap\Sigma_s\cap\bigcap_{i\ne i_0}\{g_i\le0\}.
\]
$\Sigma_s$ is a $C^2$ strictly convex hypersurface (principal curvatures in
$[c_g/\Gamma_0,\,C_g/\gamma_0]$ by \eqref{eq:reg} and the formula for the
curvature of a level set, $\kappa=(\nu^\top\Hess g_{i_0}\,\nu)/|\nabla
g_{i_0}|$ for tangent $\nu$).  We regard $G^{(i_0)}_s$ as a subset of the
$(n-1)$-manifold $\Sigma_s$, equipped with its induced Riemannian metric, and
must produce on it the data of Definition~\ref{def:config} in dimension $n-1$
with $k-1$ constraints.

\begin{lemma}[Slice is an admissible curved configuration]\label{lem:slice}
There are slice constants $c'_K\le C'_K$, $c'_g\le C'_g$, $\gamma'>0$,
$r'_0>0$ and $\delta_0>0$, depending on the fixed original configuration
and on $n,k$, such that for every $|s|\le\delta_0$ the sliced configuration
$G^{(i_0)}_s\subset\Sigma_s$ is admissible with structure constants
$(c'_K,C'_K,c'_g,C'_g,\gamma',r'_0)$, with stratum
$E^{(i_0)}_s:=E\cap\Sigma_s$ of dimension $j-1=(n-1)-1-(k-1)$.
\end{lemma}

\begin{proof}
We verify the three ingredients of Definition~\ref{def:config} for the body
$K\cap\Sigma_s$ inside the manifold $\Sigma_s$ and the constraints
$\{g_i|_{\Sigma_s}\le0\}$, $i\ne i_0$.

\emph{(a) The sliced body is strictly convex with bounded curvature.}  The body
for the lower-dimensional section lemma is the convex region cut from the convex
surface $\Sigma_s$ by $\partial K$; its boundary in $\Sigma_s$ is the
transversal intersection $\partial K\cap\Sigma_s$ of the two strictly convex
hypersurfaces $\partial K$ and $\Sigma_s$.  Its normal curvature in $\Sigma_s$
is given by the same intersection-curvature formula \eqref{eq:intersection}
below, with $\mathrm{I\!I}_{g_i}$ replaced by $\mathrm{I\!I}_K\succeq c_KI$ and
$\omega=\angle(\nu_K,\nu_{\Sigma_s})$ the dihedral angle, bounded away from
$0,\pi$ by transversality.  Since $\Sigma_s$ is a truncation surface of $K$ with
the body on the inward side, the two surfaces are mutually convex
($\cos\omega\le0$ generically), so the numerator is $\succeq c_K|X|^2$ and the
sliced body has curvature in a band $[c'_K,C'_K]$ fixed by the original
fixed-configuration data.  This yields $c'_K,C'_K$.

\emph{(b) The restricted constraints are strictly convex hypersurfaces of
$\Sigma_s$ with bounded curvature.}  This is the delicate ingredient, and we
state the mechanism, the regime where it is clean, and the regime where its
\emph{uniformity} is part of the open problem of \S\ref{sec:uniform}.  What we
need is not that the function $\tilde g_i:=g_i|_{\Sigma_s}$ be convex, but that
the truncating \emph{surface} $\Sigma'_i:=\{g_i=0\}\cap\Sigma_s$ be a strictly
convex $C^2$ hypersurface of $\Sigma_s$, curving toward the kept region, with
principal curvatures (intrinsic to $\Sigma_s$) bounded in a band $[c'_g,C'_g]$.
The surface $\Sigma'_i$ is the transversal intersection of the two convex
hypersurfaces $\{g_i=0\}$ and $\Sigma_s$.  Its normal curvature in $\Sigma_s$,
in a tangent direction $X\in T\Sigma'_i$, is given by the intersection-curvature
formula
\begin{equation}\label{eq:intersection}
\kappa_{\Sigma'_i,\Sigma_s}(X)
=\frac{\mathrm{I\!I}_{g_i}(X,X)-\cos\omega\,\mathrm{I\!I}_{\Sigma_s}(X,X)}{\sin\omega},
\qquad\omega:=\angle\bigl(n_i,\nu_{\Sigma_s}\bigr),
\end{equation}
There is only one power of \(\sin\omega\): the unit conormal of the intersection
inside \(\Sigma_s\) is \((n_i-\cos\omega\,\nu_{\Sigma_s})/\sin\omega\).
where $\mathrm{I\!I}_{g_i},\mathrm{I\!I}_{\Sigma_s}\succeq0$ are the (convex)
second fundamental forms of the two surfaces and $\omega$ is their dihedral
angle, bounded away from $0$ and $\pi$ by the transversality constant
$\gamma$ (so $\sin\omega\ge\gamma'$, $|\cos\omega|\le1-\gamma''$).  The numerator
is $\succeq c_g\,|X|^2-(1-\gamma'')C_g\,|X|^2$; this is the term that can fail to
be positive if the truncation curvature $c_g$ is small relative to $C_g$ at an
oblique angle.

\emph{Clean regime.}  If the two surfaces are \emph{mutually convex} at the
contact ($\cos\omega\le0$, normals splaying apart, the generic transversal
configuration for two truncations bounding a common region), then the numerator
is $\succeq c_g|X|^2$ and $\kappa_{\Sigma'_i,\Sigma_s}\in[c_g,(C_g+C_g)/\gamma']$:
the sliced surface is strictly convex with curvature bounded by structure
constants, giving $c'_g,C'_g$.  For our standing example (a ball truncated by
spheres whose centers lie outside, so the truncation normals splay), and more
generally whenever the truncations bound $G$ from genuinely distinct directions,
this is the operative case and the slice is unconditionally in class.

\emph{Oblique regime (uniformity caveat).}  When $\cos\omega>0$ and
$c_g<(1-\gamma'')C_g$, \eqref{eq:intersection} does not by itself certify
positivity, and the strict convexity of $\Sigma'_i$ uses the global convexity of
$\{g_i\le0\}\cap\Sigma_s$ rather than the pointwise curvature: the kept region is
still convex (an intersection of convex sets restricted to the convex surface
$\Sigma_s$), so $\Sigma'_i$ is convex toward it, but its curvature lower bound
$c'_g$ now depends on the angle $\omega$ and degrades as $\omega\to0$.  For a
\emph{fixed} configuration $\omega$ is bounded away from $0$ along the compact
stratum, so $c'_g>0$ is a genuine constant and the induction proceeds; a bound on
$c'_g$ uniform over all families with the given structure constants, as
$\omega$ ranges, is exactly Problem~\ref{prob:factor}.  The gradient bounds
$\gamma'_0,\Gamma'_0$ are the slab-projections of $\nabla g_i$, bounded by
structure constants in this clean regime and by the fixed-configuration data in
the general fixed-family regime.

\emph{(c) Transversality survives.}  The $k$ outward normals
$\{\nu_K,n_i:i\ne i_0\}$ restricted to $T\Sigma_s$ are the slab projections of
the original normals onto the tangent space of $\Sigma_s$.  Their Gram matrix is
the original Gram matrix $M$ with the row and column $i_0$ removed and the rest
projected orthogonally to $n_{i_0}$; this is the Schur complement of $M$ with
respect to its $i_0$-th entry, which is $\succeq\gamma I$ because $M\succeq\gamma
I$ and the Schur complement of a matrix $\succeq\gamma I$ along a unit direction
is again $\succeq\gamma I$ on the orthogonal complement.  Hence $\gamma'=\gamma$
works.  Finally $G^{(i_0)}_s$ contains a ball of radius $r'_0\ge c\,r_0$ in
$\Sigma_s$ for $|s|\le\delta_0$, because the inner ball $B(c^-,r_0)\subseteq G$
of the original configuration meets the slab of slices in a cap of inradius
$\ge c r_0$ inside $\Sigma_s$.
\end{proof}

\subsection{The induction and the slab bound}\label{ss:induction}

With the slice in class, the induction of \cite[\S4, Case 2]{companion11} runs
verbatim with the affine slice replaced by the curved slice $\Sigma_s$ and the
flat radius $\sqrt{1-s^2}$ replaced by the bounded geometry of
Lemma~\ref{lem:slice}.

\begin{proposition}[Cap mass, all directions]\label{prop:allcap}
For each fixed configuration satisfying Definition~\ref{def:config} there are
$b>0$ and $h_0>0$, depending on that configuration, such that for every unit $u$ with
$\tau_{\max}(u)$ attained on $E$ and every $0<h\le h_0$,
\[
\Lambda_u(h):=\mu\bigl(\{x\in G:\ \langle x,u\rangle\ge\tau_{\max}(u)-h\}\bigr)
\;\ge\;b\,h^{\frac{2n-j}2}.
\]
\end{proposition}

\begin{proof}[Proof sketch]
Induction on $k$.  The base case $k=0$ ($G$ a strictly convex body) is the
rolling-ball cap bound of \cite{companion4}: a cap of depth $h$ on a body with
principal curvatures in $[c_K,C_K]$ has the section identity
$r(h)^2=(\rho-g+h)(\rho+g-h)\ge h\rho$ with $\rho=1/C_K$, hence mass
$\ge c\,h^{(n+1)/2}=c\,h^{(2n-j)/2}$ at $j=n-1$.

For the inductive step fix $u$ with extreme on $E$.  If $u\in\mathcal U_p$ at
facet-distance $\ge\theta_0$ for a fixed threshold $\theta_0>0$ chosen from the
fixed configuration data,
Theorem~\ref{thm:cap} gives the bound directly.  Otherwise $u$ is within
$\theta_0$ of a facet of the normal cone.  Write $u=\sum_\bullet\lambda_\bullet
\nu_\bullet$ with $\lambda^\top M\lambda=1$, $M\succeq\gamma I$; then
$\max_\bullet\lambda_\bullet\ge(k+1)^{-1}\gamma^{1/2}$, so \emph{some} active
constraint $g_{i_0}$ has a multiplier bounded below, hence (by
$\beta_{i_0}=(M\lambda)_{i_0}\ge\gamma\lambda_{i_0}$) angle bounded away from
$\pi/2$: $\beta_{i_0}\ge c>0$ a structure constant.  This strongly-active
constraint is the one to slice (exactly as \cite{companion11} slices the
coordinate $i_0$ with $u_{i_0}$ bounded below in its Case~2).  For $s$ in an interval of length $\asymp h$ below the
extreme level, the slice $G^{(i_0)}_s$ is an admissible $(n-1,k-1)$ curved
configuration (Lemma~\ref{lem:slice}) and the restricted direction
$u':=P_{\Sigma_s}u/|P_{\Sigma_s}u|$ has its extreme on the sliced stratum
$E^{(i_0)}_s$, at depth $h_*\asymp h$.  By the induction hypothesis applied in
dimension $n-1$ with $k-1$ constraints (stratum dimension
$(n-1)-1-(k-1)=j-1$),
\[
\Lambda^{(n-1,k-1)}_{u'}(h_*)\;\ge\;b'\,h^{\frac{2(n-1)-(j-1)}2}
=b'\,h^{\frac{2n-j-1}2}.
\]
Integrating over the slicing interval ($\asymp h$) by the curved coarea formula
(the Jacobian $|\nabla g_{i_0}|$ is bounded above and below by \eqref{eq:reg},
so the coarea factor is a structure constant) gives
\[
\Lambda_u(h)\;\ge\;c\cdot h\cdot b'\,h^{\frac{2n-j-1}2}
\;=\;c'\,h^{\frac{2n-j}2}.
\]
The exponent telescopes exactly as in the flat case: one slice supplies a factor
$h$ (the slicing interval), and the induction supplies
$h^{(2n-j-1)/2}$, so the product is $h^{1+(2n-j-1)/2}=h^{(2n-j)/2}$, with the
codimension count $j\to j-1$ matching $(n,k)\to(n-1,k-1)$.
\end{proof}

\begin{proof}[Proof of Theorem~\ref{thm:slab}]
This is the anchor--unimodality conversion of \cite{companion4,companion11},
which uses only: (i) the cap-mass bound of Proposition~\ref{prop:allcap} at both
$u$-extremes; (ii) the unimodality of the level-section area $A_u$, which holds
by Brunn--Minkowski because $G$ is convex; and (iii) the inner ball
$B(c^-,r_0)\subseteq G$ providing a bulk anchor of value $a_1(n,k)>0$ at a depth
$h^\dagger\ge c(k)>0$.  Define $A_u(h):=\vol_{n-1}(G\cap\{\langle x,u\rangle
=\tau_{\max}-h\})$.  Below the mode of $A_u$, $A_u(h)\ge\Lambda_u(h)/h\ge
b\,h^{(2n-j-1)/2}$ by Proposition~\ref{prop:allcap}; above the mode,
unimodality gives $A_u(h)\ge A_u(h^\dagger)\ge a_1$.  Anchoring at depth $L/4$
from each of the $u$- and $(-u)$-extremes and integrating over the middle of the
$\tau$-range (whose complement has length $\le L/2$, while the range has length
$\ge2r_0\ge2L$) yields
$\vol(G\cap\{|\langle y-x,u\rangle|\le L\})\ge a_1 L^{(2n-j)/2}$.
\end{proof}

\begin{proof}[Proof of Corollary~\ref{cor:law}]
\emph{Upper bound.}  Insert Theorem~\ref{thm:slab} into the proof skeleton of
\cite{companion2,companion4}: the population dichotomy holds with the inner-ball
radius $r_0$; the VC relative deviation bounds apply to slabs cut by $k$ fixed
$C^2$ convex constraints, whose graphs $\{g_i\le0\}$ together with two parallel
slab walls form a family of sets of finite VC dimension (the class
$\{g_i\le t\}$ is a translate of a fixed graph, hence a VC class, and finite
intersections multiply the dimension); thus a sample point in $S_w$ yields
$\mu(S_{\kappa w})\ge32u_N$ for all $w\ge\ell_N:=C(u_N)^{2/(2n-j)}\asymp
(\log N/N)^{2/(2n-j)}$, and the UDC holds with probability $1-1/N$.

\emph{Lower bound.}  The second-moment scheme of
\cite{companion4,companion11} applies with the curved stratum caps of
Theorem~\ref{thm:cap} in place of the flat ones.  The cap directions are the
stratum directions $u_p\in\mathcal U_p$ at a fixed facet-distance $\theta_0$
(any fixed interior direction of the normal cone), so Theorem~\ref{thm:cap}
gives two-sided cap masses $\asymp h^{(2n-j)/2}$ with $\theta_0$-dependent but
$N$-independent constants.  The footprint of a depth-$H$ cap is $\asymp\sqrt H$
(the $y$-ball radius), so a maximal $2c\sqrt{H_N}$-separated family on the
$j$-manifold $E$ has $M\ge c H_N^{-j/2}\ge c'(N/\tau_N)^{j/(2n-j)}$ members with
pairwise disjoint outer caps.  The first/second moment computation is identical
to \cite{companion11}: with $\tau_N=\varepsilon\log N$ chosen so that the
per-cap gap probability is $N^{-2q\varepsilon}$ with
$q=(4\kappa^2)^{(2n-j)/2}b_2/b_1$, one gets $\E X\to\infty$ and
$\PP(X=0)=o(1)$ by Chebyshev, forcing $\ell\ge c(\log N/N)^{2/(2n-j)}$.
\end{proof}

\section{What uniformity is and is not obtained}\label{sec:uniform}

We now state precisely the scope of Corollary~\ref{cor:law} and the residual
open problem.

\subsection{Obtained: fixed family, uniform over the stratum and compact subcones}

For a \emph{fixed} admissible configuration (a fixed body $K$ and fixed $C^2$
strictly convex $g_1,\dots,g_k$), Corollary~\ref{cor:law} gives the exact law
$\Theta((\log N/N)^{2/(2n-j)})$.  The constants in
Theorems~\ref{thm:cap}--\ref{thm:slab} are uniform:
\begin{itemize}
\item \emph{over the stratum} $p\in E$: the structure constants
$(c_K,C_K,c_g,C_g,\gamma_0,\Gamma_0,\gamma,r_0)$ and the $C^2$ Taylor remainders
are uniform on the compact manifold $E$;
\item \emph{over directions} $u$ in any fixed compact subcone of the normal cone
(facet-distance $\ge\theta_0$): the cap exponent is exactly $(2n-j)/2$, with the
constant determined by the fixed configuration and \(\theta_0\), for all depths
$h\le h_1\theta_0^2$.  The upper bound constant $b_2$ is uniform over the cone,
while the lower bound constant $b_1(\theta_0)$ degrades as the subcone
approaches the boundary; this is the loss coming from the pinned lengths
$s_i\le Ch/(\beta_*\theta_0)$.
\end{itemize}
This is the regime needed for both directions of the law: the slab-volume bound
(Theorem~\ref{thm:slab}) is a statement for \emph{all} $u$, obtained by the
slicing induction which always slices the most binding constraint and so always
operates at definite angle; and the lower-bound caps may be taken at a single
fixed interior direction $\theta_0=\frac12\operatorname{diam}\mathcal U_p$, say.

\subsection{The exponent is curvature-robust but direction-sensitive}

The exponent $(2n-j)/2$ does not depend on the magnitudes of the curvatures: for
a fixed active set and fixed interior direction the slopes $m_i$ are fixed
constants, so $h\to0$ keeps every constraint steep ($m_i\le C_*^{-1/2}h^{-1/2}$
automatically once $h$ is small), and each pinned coordinate contributes a clean
factor $h$ regardless of $Q_i$.  High truncation curvature only inflates the
constant (through $\sqrt{h/c_g}$ in the shallow regime and the joint form in
Remark~\ref{rem:factor}); it never changes the exponent.

The exponent \emph{does} change as $u$ approaches the boundary of the normal
cone $\mathcal U_p$, and this is correct rather than a defect.  Two equivalent
descriptions of the boundary: a multiplier vanishes ($\lambda_{i_0}\to0$, the
$i_0$-th constraint becomes first-order inactive at the maximizer), or
equivalently the slab tilts toward tangency with surface $i_0$
($\beta_{i_0}=\langle n_{i_0},u\rangle\to0$, slope $m_{i_0}\to\infty$), so that
$p$ ceases to be the strict $\tau$-maximum and the extreme spreads onto a
\emph{larger} stratum $E'$ of dimension $j+1$.  There the cap is governed by the
$(2n-j')/2$ law with $j'=j+1$: a \emph{larger} cap.  The natural single formula
across the closed cone is the \emph{trichotomy interpolant}: with the per-surface
pinned length
\[
\lambda_i(h):=\min\Bigl(\frac{h}{m_i},\ \sqrt{h/c_g}\Bigr),
\qquad m_i=\tan\angle(n_i,u),
\]
one expects
\begin{equation}\label{eq:interp}
\mu(C(h))\;\asymp\;h^{1+j/2}\,\prod_{i=1}^k\lambda_i(h),
\end{equation}
which interpolates: when all $m_i$ are $O(1)$, $\lambda_i(h)\asymp h$ and
\eqref{eq:interp} is $\asymp h^{1+j/2}\cdot h^k=h^{(2n-j)/2}$; when $m$ of the
slopes exceed $h^{-1/2}$ (those $\beta_i\lesssim\sqrt h$), the curvature cap
gives $\lambda_i(h)\asymp\sqrt h$ and the exponent rises to $(2n-(j+m))/2$,
exactly matching the larger stratum that has become active.  The crossover is at
$m_i\asymp h^{-1/2}$, i.e.\ $\beta_i\asymp\sqrt h$, the same $\sqrt h$ angular
threshold as the simplicial trichotomy $\mathcal A/\mathcal Z$ of
\cite{companion11} (there $\lambda_i=\min(h/(8ku_i),\sqrt h/8)$ with
$u_i=\beta_i$); curvature does not move the threshold, since the shallow length
$\sqrt{h/c_g}$ and the steep length $h/m_i$ are sandwiched between the flat-case
values with constants from $[c_g,C_g]$.  We obtain \eqref{eq:interp} as a valid
\emph{upper} bound for all $u$ (the per-surface lengths $\lambda_i(h)$ each cap
the $i$-th coordinate independently); the matching lower bound is
Theorem~\ref{thm:cap} away from $\partial\mathcal U_p$, and is open in a
shrinking neighborhood of the multiply-tangent directions, where the pinned
block can degenerate jointly even though each $\lambda_i$ is controlled.

\subsection{Open: uniformity over curvature-bounded families}

The hypotheses we cannot remove --- and which constitute the residual open
problem --- concern uniformity \emph{across a family} of configurations rather
than for a single fixed one.  Two distinct gaps:

\begin{problem}[Uniform constant near multiply-tangent directions]\label{prob:uniform}
Let $\mathcal F(c_K,C_K,c_g,C_g,\gamma_0,\Gamma_0,\gamma,r_0)$ be the class of
\emph{all} admissible configurations with the given structure constants.  Is the
slab-volume bound of Theorem~\ref{thm:slab} uniform over $\mathcal F$ --- i.e.\
are $a_1,L_0$ functions of the structure constants alone, with no dependence on
the individual configuration?  Equivalently, in \eqref{eq:interp}, is the
\emph{lower} bound uniform as $u$ approaches a codimension-$\ge2$ face of the
normal cone, where two or more $\beta_i\to0$ simultaneously?
\end{problem}

\noindent The obstruction is genuine.  When a single $\beta_{i_0}\to0$, the
slicing induction handles it (slice on $g_{i_0}$).  When two angles
$\beta_{i_0},\beta_{i_1}\to0$ together, the slab projections $P_un_{i_0},
P_un_{i_1}$ of the two near-tangent normals can themselves become nearly
parallel, so the dual frame \eqref{eq:dual} blows up, $|f_i|\to\infty$, and the
joint Gram of the pinned block degrades.  The trichotomy bound \eqref{eq:interp}
remains a valid \emph{upper} bound there (it only uses the diagonal $\beta_i$),
but the matching \emph{lower} bound requires the joint Gram of
$\{P_un_i:\beta_i\ \text{small}\}$ to stay bounded below, which is an extra
non-degeneracy hypothesis not implied by the per-surface structure constants.
For a fixed configuration this is automatic (the finite set of multiply-tangent
directions has its own positive Gram lower bound, absorbed into the
configuration's constants); uniformly over $\mathcal F$ it is open, because the
Gram lower bound can be made to deteriorate within the class while keeping the
structure constants fixed.

There is, however, a clean conditional positive statement.  It isolates exactly
the extra geometry that the proof needs beyond the per-surface structure
constants.

\begin{definition}[Resolved multi-tangent geometry]\label{def:resolved}
Fix structure constants and numbers \(\Gamma,c_{\rm sl}>0\).  A subfamily
\(\mathcal F_{\Gamma,c_{\rm sl}}\subseteq
\mathcal F(c_K,C_K,c_g,C_g,\gamma_0,\Gamma_0,\gamma,r_0)\) is called
\emph{\((\Gamma,c_{\rm sl})\)-resolved} if the following two properties hold
for every configuration in the subfamily and for every iterated one-constraint
slice that can occur in the induction of Lemma~\ref{lem:slice}.
\begin{enumerate}
\item[(R1)] Every sliced truncating surface remains strictly convex with lower
principal curvature at least \(c_{\rm sl}\) and with upper curvature and
gradient bounds controlled by the original structure constants and
\(c_{\rm sl}^{-1}\).
\item[(R2)] Let \(p\) be a point of the active stratum, let \(u\) lie in the
closed normal cone at \(p\), and let \(P_u\) denote orthogonal projection onto
\(u^\perp\).  For every subset \(I\) of active curved truncating normals for
which \(P_un_i(p)\ne0\), the normalized projected vectors
\[
e_i(u,p):=\frac{P_un_i(p)}{|P_un_i(p)|},\qquad i\in I,
\]
have Gram matrix bounded below by \(\Gamma I\).
\end{enumerate}
Condition (R1) is automatic in the clean slicing regime of
Lemma~\ref{lem:slice}; condition (R2) is automatic for a fixed configuration
but is not forced uniformly by the structure constants alone.
\end{definition}

\begin{proposition}[Uniform curved law under resolved multi-tangent geometry]
\label{prop:resolved-uniform}
On every \((\Gamma,c_{\rm sl})\)-resolved subfamily
\(\mathcal F_{\Gamma,c_{\rm sl}}\), the slab bound of
Theorem~\ref{thm:slab} is uniform: there are \(a_*,L_*>0\), depending only on
the structure constants, \(n,k,\Gamma\), and \(c_{\rm sl}\), such that every
configuration in \(\mathcal F_{\Gamma,c_{\rm sl}}\) satisfies
\[
\vol\Bigl(G\cap\{y:\ |\langle y-x,u\rangle|\le L\}\Bigr)
\ge a_*\,L^{\frac{2n-j}2}
\]
for every \(x\in G\), every unit \(u\), and every \(0<L\le L_*\).
Consequently the interpolation law of Corollary~\ref{cor:law} holds uniformly
over \(\mathcal F_{\Gamma,c_{\rm sl}}\), with UDC constants depending only on
the displayed parameters and not on the individual curved configuration.
\end{proposition}

\begin{proof}
Repeat the proof of Proposition~\ref{prop:allcap} and
Theorem~\ref{thm:slab}, tracking constants.  Away from a face of the normal
cone, the lower witness box in Theorem~\ref{thm:cap} only uses bounded
Taylor remainders, the curvature bands, and the bounded inverse of the Gram
matrix of the projected active normals.  Condition (R2) supplies that inverse
bound uniformly, so the dual frame \eqref{eq:dual} has
\(|f_i|\le C(\Gamma,k)\) throughout the resolved subfamily.

Near a face, the induction slices on a strongly active constraint exactly as in
the proof of Proposition~\ref{prop:allcap}.  Condition (R1) says that the sliced
configuration remains in the same quantitative class, while the Schur-complement
transversality argument in Lemma~\ref{lem:slice} gives the same lower Gram bound
for the remaining unsliced normals.  Condition (R2), imposed on all iterated
slices, prevents the only remaining loss: simultaneous projected-normal
degeneration in the pinned block.  Hence the inductive constants in the cap
lower bound are functions only of \(n,k\), the structure constants,
\(\Gamma\), and \(c_{\rm sl}\).  The coarea interval has length \(\asymp h\)
with constants controlled by the gradient bounds, so the exponent telescopes as
before and yields a uniform version of Proposition~\ref{prop:allcap}.

The conversion from uniform cap mass to the uniform slab estimate is the same
anchor--unimodality argument used in Theorem~\ref{thm:slab}; the inner-ball
radius is part of the structure constants.  Finally the empirical UDC proof
skeleton of the companion papers uses only this slab constant, the fixed VC
dimension of the family, and the density bounds, so the high-probability
interpolation law is uniform on the resolved subfamily.
\end{proof}

\begin{corollary}[Finite curved atlases are uniform]
\label{cor:finite-atlas}
Let \(\mathcal F_0=\{(K^a,g^a_1,\ldots,g^a_k):1\le a\le M\}\) be a finite
collection of admissible curved configurations with the same \(n,k\).  After
enlarging the structure constants to dominate the finitely many members, there
are \(a_0,L_0>0\), depending on the finite atlas \(\mathcal F_0\), such that
every configuration in \(\mathcal F_0\) satisfies
\[
\vol\Bigl(G^a\cap\{y:\ |\langle y-x,u\rangle|\le L\}\Bigr)
\ge a_0\,L^{\frac{2n-j}2}
\]
for every \(a\), every \(x\in G^a\), every unit \(u\), and every
\(0<L\le L_0\).  Consequently, for any common density bounds, the interpolation
law of Corollary~\ref{cor:law} holds uniformly over the whole finite atlas.
\end{corollary}

\begin{proof}
For each fixed member of the atlas, the discussion following
Definition~\ref{def:resolved} gives constants
\(\Gamma_a,c_{{\rm sl},a}>0\) for which that single configuration is
\((\Gamma_a,c_{{\rm sl},a})\)-resolved: compactness and fixed-configuration
admissibility give positive Gram minima for the projected normals and positive
curvature floors for the slices.  Since the atlas is finite,
\[
  \Gamma_0:=\min_a\Gamma_a>0,\qquad
  c_0:=\min_a c_{{\rm sl},a}>0 .
\]
Thus \(\mathcal F_0\) is \((\Gamma_0,c_0)\)-resolved after using common
structure constants for all members.  Proposition~\ref{prop:resolved-uniform}
gives the displayed slab estimate and the uniform interpolation law.
\end{proof}

\begin{corollary}[Compact resolved curved atlases are uniform]
\label{cor:compact-atlas}
Let \(T\) be a compact parameter space, and let
\[
  \mathcal F_T=\{(K^\tau,g^\tau_1,\ldots,g^\tau_k):\tau\in T\}
\]
be a family of admissible curved configurations with the same \(n,k\) and
common structure constants.  Suppose there are continuous functions
\[
  \Gamma:T\to(0,\infty),\qquad c_{\rm sl}:T\to(0,\infty),
\]
such that the configuration indexed by \(\tau\) is
\((\Gamma(\tau),c_{\rm sl}(\tau))\)-resolved.  Then there are \(a_0,L_0>0\),
depending on the compact atlas \(\mathcal F_T\), such that every
\(\tau\in T\) satisfies
\[
\vol\Bigl(G^\tau\cap\{y:\ |\langle y-x,u\rangle|\le L\}\Bigr)
\ge a_0\,L^{\frac{2n-j}2}
\]
for every \(x\in G^\tau\), every unit \(u\), and every \(0<L\le L_0\).
Consequently, for any common density bounds, the interpolation law of
Corollary~\ref{cor:law} holds uniformly over the whole compact atlas.
\end{corollary}

\begin{proof}
By compactness and continuity,
\[
  \Gamma_0:=\min_{\tau\in T}\Gamma(\tau)>0,\qquad
  c_0:=\min_{\tau\in T}c_{\rm sl}(\tau)>0 .
\]
Thus the whole atlas is \((\Gamma_0,c_0)\)-resolved.  Apply
Proposition~\ref{prop:resolved-uniform}.
\end{proof}

\begin{corollary}[\(C^2\)-neighborhoods of resolved configurations]
\label{cor:resolved-neighborhood}
Let
\[
  \{(K^\tau,g^\tau_1,\ldots,g^\tau_k):\tau\in T\}
\]
be a parametrized family of admissible curved configurations with common
structure constants, and suppose the data vary continuously in \(C^2\) in the
following local sense: after choosing finitely many graph charts for the active
stratum and for every iterated one-constraint slice occurring in
Lemma~\ref{lem:slice}, the graph functions, their first derivatives, and their
Hessians depend continuously on \(\tau\) on compact chart domains.  If the
configuration at \(\tau_0\) is \((\Gamma_0,c_0)\)-resolved, then there is a
neighborhood \(U\) of \(\tau_0\) such that every configuration with
\(\tau\in U\) is \((\Gamma_0/2,c_0/2)\)-resolved, after possibly restricting
the slice-depth parameter in Lemma~\ref{lem:slice} by a factor depending only
on the common structure constants and on the chosen charts.

Consequently, on \(U\) the slab bound of Theorem~\ref{thm:slab} and the
interpolation law of Corollary~\ref{cor:law} hold uniformly, with constants
depending only on the common structure constants, \(n,k,\Gamma_0\), and \(c_0\).
\end{corollary}

\begin{proof}
For the fixed configuration \(\tau_0\), the quantities in (R1) and (R2) are
positive minima of continuous functions over compact incidence spaces: the
active stratum is compact, there are only finitely many orders in which the
induction can slice the \(k\) constraints, the permitted slice parameters lie
in a compact interval after the standard depth restriction, the unit normal
directions and pinned subsets range over compact sets, and the list of active
subsets is finite.  By hypothesis these functions vary continuously in the
chosen \(C^2\) charts.  Hence their minima are lower semicontinuous at
\(\tau_0\).  After shrinking the parameter neighborhood \(U\), all sliced
curvature floors stay above \(c_0/2\), and all projected-normal Gram minima in
(R2) stay above \(\Gamma_0/2\).  Upper curvature, gradient, and transversality
bounds remain controlled because the family has common structure constants and
the charts vary \(C^2\)-continuously.  Thus the nearby configurations are
\((\Gamma_0/2,c_0/2)\)-resolved.  Proposition~\ref{prop:resolved-uniform}
then gives the uniform slab and UDC conclusions.
\end{proof}

\begin{corollary}[Compact \(C^2\)-continuous resolved atlases]
\label{cor:c2-compact-atlas}
Let \(T\) be compact, and let
\[
  \mathcal F_T=\{(K^\tau,g^\tau_1,\ldots,g^\tau_k):\tau\in T\}
\]
be a \(C^2\)-continuous atlas in the sense of
Corollary~\ref{cor:resolved-neighborhood}, with common structure constants.
Assume that each individual configuration in the atlas is resolved in the
sense of Definition~\ref{def:resolved}.  Then there are \(a_0,L_0>0\) such
that every \(\tau\in T\) satisfies
\[
\vol\Bigl(G^\tau\cap\{y:\ |\langle y-x,u\rangle|\le L\}\Bigr)
\ge a_0\,L^{\frac{2n-j}2}
\]
for every \(x\in G^\tau\), every unit \(u\), and every \(0<L\le L_0\).
Consequently, for any common density bounds, the interpolation law of
Corollary~\ref{cor:law} holds uniformly over the whole compact atlas.
\end{corollary}

\begin{proof}
For each \(\tau\in T\), choose resolved constants
\((\Gamma_\tau,c_\tau)\) for that single configuration.  By
Corollary~\ref{cor:resolved-neighborhood}, there is a neighborhood
\(U_\tau\) on which the atlas is
\((\Gamma_\tau/2,c_\tau/2)\)-resolved.  Compactness gives a finite subcover
\(U_{\tau_1},\ldots,U_{\tau_M}\).  Set
\[
  \Gamma_*:=\frac12\min_{1\le a\le M}\Gamma_{\tau_a}>0,\qquad
  c_*:=\frac12\min_{1\le a\le M}c_{\tau_a}>0 .
\]
Then the whole atlas is \((\Gamma_*,c_*)\)-resolved after using the common
structure constants.  Proposition~\ref{prop:resolved-uniform} gives the
uniform slab estimate and the uniform interpolation law.
\end{proof}

\begin{corollary}[Exhaustive resolved curved atlases]
\label{cor:exhaustive-atlas}
Let \(T=\bigcup_{q\ge1}T_q\), where \(T_q\subset T_{q+1}\) are compact
parameter spaces, and let
\[
  \mathcal F_T=\{(K^\tau,g^\tau_1,\ldots,g^\tau_k):\tau\in T\}
\]
be a family of admissible curved configurations with the same \(n,k\) and common
structure constants.  Suppose there are functions
\[
  \Gamma:T\to(0,\infty),\qquad c_{\rm sl}:T\to(0,\infty),
\]
continuous on each \(T_q\), such that the configuration indexed by \(\tau\) is
\((\Gamma(\tau),c_{\rm sl}(\tau))\)-resolved.  Then, for each \(q\), there are
constants \(a_q,L_q>0\) such that
\[
\vol\Bigl(G^\tau\cap\{y:\ |\langle y-x,u\rangle|\le L\}\Bigr)
\ge a_q\,L^{\frac{2n-j}2}
\]
for every \(\tau\in T_q\), every \(x\in G^\tau\), every unit \(u\), and every
\(0<L\le L_q\).  Consequently, for any common density bounds, the interpolation
law of Corollary~\ref{cor:law} holds uniformly on each compact stage \(T_q\).

If, in addition,
\[
  \inf_{\tau\in T}\Gamma(\tau)>0,\qquad
  \inf_{\tau\in T}c_{\rm sl}(\tau)>0,
\]
then \(a_q,L_q\) may be chosen independently of \(q\), and the interpolation law
is uniform over the whole atlas \(\mathcal F_T\).
\end{corollary}

\begin{proof}
For each fixed \(q\), Corollary~\ref{cor:compact-atlas} applies to the compact
subatlas over \(T_q\), giving constants \(a_q,L_q\).  If the two displayed
infima over \(T\) are positive, then the whole family lies in one resolved class
with
\[
  \Gamma_0=\inf_{\tau\in T}\Gamma(\tau),\qquad
  c_0=\inf_{\tau\in T}c_{\rm sl}(\tau).
\]
Proposition~\ref{prop:resolved-uniform} gives one pair of constants for all
\(\tau\in T\).
\end{proof}

\begin{corollary}[Uniform one-cut curved law under splaying]
\label{cor:onecut-uniform}
Fix structure constants as in Definition~\ref{def:config} and fix
\(\chi>0\).  Let \(\mathcal F_{1,\chi}\) be the class of admissible curved
configurations with \(k=1\) such that, for every \(p\in E\),
\[
  \langle \nu_K(p),n_1(p)\rangle\le -\chi .
\]
Then the slab bound of Theorem~\ref{thm:slab} is uniform on
\(\mathcal F_{1,\chi}\): there are \(a_\chi,L_\chi>0\), depending only on
\(n\), the structure constants, and \(\chi\), such that every configuration in
\(\mathcal F_{1,\chi}\) satisfies
\[
\vol\Bigl(G\cap\{y:\ |\langle y-x,u\rangle|\le L\}\Bigr)
\ge a_\chi\,L^{\frac{n+2}2}
\]
for every \(x\in G\), every unit \(u\), and every \(0<L\le L_\chi\).
Consequently the interpolation law is uniform on \(\mathcal F_{1,\chi}\):
\[
  \ell_N\asymp\Bigl(\frac{\log N}{N}\Bigr)^{2/(n+2)}
\]
with constants depending only on \(n\), the density bounds, the structure
constants, and \(\chi\).
\end{corollary}

\begin{proof}
For \(k=1\) there is no simultaneous multi-tangent block of curved truncating
normals.  Thus condition (R2) in Definition~\ref{def:resolved} is automatic
with \(\Gamma=1\), since every relevant projected-normal family has at most one
unit vector.

It remains to check that the one possible slice has uniform geometry.  Slice by
\(g_1=s\).  After this slice there are no remaining truncating surfaces, so the
only lower-dimensional body is \(K\cap\Sigma_s\subset\Sigma_s\), where
\(\Sigma_s=\{g_1=s\}\).  The splaying condition and the \(C^2\) bounds imply,
after choosing \(|s|\le\delta_\chi\) with \(\delta_\chi\) depending only on the
displayed parameters, that the dihedral angle between \(\partial K\) and
\(\Sigma_s\) still satisfies
\[
  \cos\omega\le-\chi/2,\qquad \sin\omega\ge c(\gamma,\chi)>0 .
\]
The intersection-curvature formula \eqref{eq:intersection}, with
\(\mathrm{I\!I}_{g_i}\) replaced by \(\mathrm{I\!I}_K\), then gives
\[
  \kappa_{\partial K\cap\Sigma_s,\Sigma_s}(X)
  =
  \frac{\mathrm{I\!I}_K(X,X)-\cos\omega\,\mathrm{I\!I}_{\Sigma_s}(X,X)}
       {\sin\omega}
  \ge c_\chi |X|^2
\]
and also an upper curvature bound \(C_\chi |X|^2\), with constants depending
only on the structure constants and \(\chi\).  The sliced inradius and
transversality bounds are inherited from the original inner-ball and Gram
assumptions, as in Lemma~\ref{lem:slice}.  Hence the family
\(\mathcal F_{1,\chi}\) is \((1,c_\chi)\)-resolved in the sense of
Definition~\ref{def:resolved}.  Proposition~\ref{prop:resolved-uniform} gives
the uniform slab estimate and the uniform interpolation law.  Since
\(j=n-1-k=n-2\), the exponent is \(2/(2n-j)=2/(n+2)\).
\end{proof}

\begin{corollary}[Uniform one-cut curved law under curvature dominance]
\label{cor:onecut-curvdom}
Fix structure constants as in Definition~\ref{def:config} and fix
\(\chi>0\).  Let \(\mathcal F^{\rm dom}_{1,\chi}\) be the class of admissible
curved configurations with \(k=1\) such that, for every \(p\in E\),
\[
  c_K-\max\{\langle\nu_K(p),n_1(p)\rangle,0\}\,\frac{C_g}{\gamma_0}\ge\chi .
\]
Then the slab bound of Theorem~\ref{thm:slab} is uniform on
\(\mathcal F^{\rm dom}_{1,\chi}\): there are \(a_{\rm dom},L_{\rm dom}>0\),
depending only on \(n\), the structure constants, and \(\chi\), such that every
configuration in \(\mathcal F^{\rm dom}_{1,\chi}\) satisfies
\[
\vol\Bigl(G\cap\{y:\ |\langle y-x,u\rangle|\le L\}\Bigr)
\ge a_{\rm dom}\,L^{\frac{n+2}2}
\]
for every \(x\in G\), every unit \(u\), and every \(0<L\le L_{\rm dom}\).
Consequently the interpolation law is uniform on
\(\mathcal F^{\rm dom}_{1,\chi}\):
\[
  \ell_N\asymp\Bigl(\frac{\log N}{N}\Bigr)^{2/(n+2)}.
\]
In particular, if \(c_K>C_g/\gamma_0\), then every one-cut admissible family
with the given structure constants is covered, with
\(\chi=c_K-C_g/\gamma_0\).
\end{corollary}

\begin{proof}
As in Corollary~\ref{cor:onecut-uniform}, condition (R2) of
Definition~\ref{def:resolved} is automatic for \(k=1\), because there is at
most one projected truncating normal in any pinned block.  We only need a
uniform lower curvature bound for the single slice \(K\cap\Sigma_s\subset
\Sigma_s\), where \(\Sigma_s=\{g_1=s\}\).

The level-surface curvature bound in the proof of Lemma~\ref{lem:slice} gives
\(\mathrm{I\!I}_{\Sigma_s}\preceq (C_g/\gamma_0)I\).  The margin assumption and
the \(C^2\) bounds imply that, after restricting to \(|s|\le\delta_\chi\) with
\(\delta_\chi\) depending only on the displayed parameters,
\[
  c_K-\max\{\cos\omega,0\}\,\frac{C_g}{\gamma_0}\ge\chi/2,
  \qquad \sin\omega\ge c(\gamma)>0,
\]
for the dihedral angle \(\omega\) between \(\partial K\) and \(\Sigma_s\).
Therefore the intersection-curvature formula \eqref{eq:intersection}, again
with \(\mathrm{I\!I}_{g_i}\) replaced by \(\mathrm{I\!I}_K\), yields
\[
  \kappa_{\partial K\cap\Sigma_s,\Sigma_s}(X)
  =
  \frac{\mathrm{I\!I}_K(X,X)-\cos\omega\,\mathrm{I\!I}_{\Sigma_s}(X,X)}
       {\sin\omega}
  \ge c_{\rm dom}|X|^2,
\]
with \(c_{\rm dom}>0\) depending only on the structure constants and \(\chi\).
The corresponding upper curvature, inradius, and transversality bounds are
controlled in the same way as in Lemma~\ref{lem:slice}.  Hence
\(\mathcal F^{\rm dom}_{1,\chi}\) is \((1,c_{\rm dom})\)-resolved, and
Proposition~\ref{prop:resolved-uniform} gives the uniform slab estimate and the
uniform interpolation law.  The final assertion follows because
\(\max\{\langle\nu_K,n_1\rangle,0\}\le1\).
\end{proof}

\begin{corollary}[Uniform one-cut law under nonobtuse contact]
\label{cor:onecut-nonobtuse}
Fix structure constants as in Definition~\ref{def:config}.  Let
\(\mathcal F^{\rm no}_{1}\) be the class of admissible curved configurations
with \(k=1\) such that, for every \(p\in E\),
\[
  \langle\nu_K(p),n_1(p)\rangle\le0 .
\]
Then the slab bound of Theorem~\ref{thm:slab} is uniform on
\(\mathcal F^{\rm no}_{1}\): there are \(a_{\rm no},L_{\rm no}>0\), depending
only on \(n\) and the structure constants, such that every configuration in
\(\mathcal F^{\rm no}_{1}\) satisfies
\[
\vol\Bigl(G\cap\{y:\ |\langle y-x,u\rangle|\le L\}\Bigr)
\ge a_{\rm no}\,L^{\frac{n+2}2}
\]
for every \(x\in G\), every unit \(u\), and every \(0<L\le L_{\rm no}\).
Consequently the interpolation law is uniform on \(\mathcal F^{\rm no}_{1}\):
\[
  \ell_N\asymp\Bigl(\frac{\log N}{N}\Bigr)^{2/(n+2)} .
\]
\end{corollary}

\begin{proof}
For every configuration in \(\mathcal F^{\rm no}_{1}\),
\[
  c_K-\max\{\langle\nu_K(p),n_1(p)\rangle,0\}\frac{C_g}{\gamma_0}
  \ge c_K
  \qquad(p\in E).
\]
Thus the class is contained in
\(\mathcal F^{\rm dom}_{1,c_K/2}\).  Corollary~\ref{cor:onecut-curvdom}
gives the uniform slab estimate and interpolation law.
\end{proof}

\begin{problem}[Separation of curvatures]\label{prob:factor}
In the simplicial model the cap constant factored through the stratum curvature
$Q$ and the truncation angles separately.  Quantify the failure for curved data:
find the sharp dependence of $b_1,b_2$ on the joint second fundamental form
$Q_K|_V+\sum_iQ_i|_V$ versus the separate forms, and determine whether a
factorized lower bound
$b_1\ge c\,(\det Q_K|_V)^{-1/2}\prod_i\phi(\text{angle}_i)$ holds with $\phi$
independent of the cross-curvature.  Remark~\ref{rem:factor} shows the cross
blocks are lower order at the cap scale, so a factorized bound holds \emph{up to
a multiplicative constant depending on the condition number of the joint form};
whether that constant can be removed --- giving a clean product formula as in the
flat case --- is open.
\end{problem}

We summarize the honest status.  \emph{Fixed smooth strictly convex family, in
general position, transversal along a $j$-stratum: the interpolation law
$\Theta((\log N/N)^{2/(2n-j)})$ holds, with the same exponent and logarithm as
the flat case, and constants determined by the configuration's curvature and
transversality data} (Corollary~\ref{cor:law}).  Away from multiply-tangent
directions the estimates are controlled by the stated nondegeneracy constants;
Corollary~\ref{cor:onecut-uniform} gives a fully uniform one-curved-cut theorem
under quantitative splaying, and Corollary~\ref{cor:onecut-curvdom} gives the
same uniform law in the curvature-dominated oblique one-cut regime, with
Corollary~\ref{cor:onecut-nonobtuse} recording the nonobtuse-contact
specialization.  The atlas
corollaries, Corollaries~\ref{cor:finite-atlas} and \ref{cor:compact-atlas},
cover any finite menu of fixed curved geometries, and any compact atlas whose
resolved constants vary continuously and stay positive; the local stability and
finite-subcover form in Corollaries~\ref{cor:resolved-neighborhood} and
\ref{cor:c2-compact-atlas} show that resolved \(C^2\)-continuous compact
atlases are uniform even when one checks the resolved geometry point by point.
The remaining curvature-band-uniform slab/lower constants are open only where
these one-cut slice-curvature criteria fail, where no compact \(C^2\)-resolved
atlas controls the degeneration, or in the shrinking multi-tangent regime
(Problems~\ref{prob:uniform}--\ref{prob:factor}).
As a fixed-family assertion, the exponent is settled for every first-order
active stratum: it is $(2n-j)/2$ for that stratum, and the trichotomy bound
\eqref{eq:interp} is the fixed-configuration interpolant across strata.

\section{Concluding remarks}\label{sec:concl}

The curved case isolates exactly which features of the simplicial proof are
essential and which are conveniences.  The orthogonality $\langle
w,\tilde u\rangle=0$ of the orthogonal model and the biorthogonal frame
$\langle f_i,a_{i'}\rangle=\delta_{ii'}$ of the general-position simplicial
proof are conveniences: both are devices to decouple the $k$ constraints in the
level functional, and both are replaced here by a single completing-the-square
in the slab coordinates, valid because the cross blocks of the joint second
fundamental form are lower order ($O(h)$) at the cap scale ($\sqrt h$).  The
essential feature is the \emph{linear pinning}: each transversal constraint
contributes a full factor $h$ through its first-order part, and this survives
curvature because the strictly convex correction is quadratic, hence $O(h^2)$ in
the pinned variable.  The only place curvature could change the exponent is the
near-tangent regime $\beta_i\lesssim\sqrt h$, and there it does not change the
exponent of the relevant stratum --- it relabels which stratum is active, with
the trichotomy threshold $\sqrt h$ identical to the flat case.

Within the stability theory of \cite{BJS,ShmaloUDC} the conclusion is the same
as for flat truncations for each fixed smooth activation geometry: the finest
resolution at which i.i.d.\ data certify the doubling condition is set by the
most isolated extremal sample point of the lowest transversal stratum of the
truncation geometry, with the exponent ladder $\frac2{n+1},\dots,\frac1n$
governed by the stratum dimension and \emph{not} by the curvature of the
truncating surfaces.  The remaining quantitative question --- a curvature-band-
and transversality-band-uniform lower/slab constant near the multiply-tangent
directions --- is the
honest residue, and is the natural next target.

\subsection*{Acknowledgments}
The research presented in this paper was supported by the European Research
Council (ERC) under the European Union's Horizon Europe research and innovation
programme (grant agreement No.\ 101041711), by the Israel Science Foundation
(grant number 2258/19) and by the Israel Science Foundation (ISF Grant
4101/25).

\subsection*{Declaration on generative AI}
During the preparation of this work the author used generative AI tools to
accelerate drafting and revision.  All mathematical claims, proofs, and
bibliographic information were subsequently reviewed and edited by the author,
who takes full responsibility for the content of the manuscript.

\begin{thebibliography}{9}

\bibitem{BJS}
L.~Berlyand, P.-E.~Jabin and C.~A.~Safsten,
\emph{Stability for the training of deep neural networks and other classifiers},
Math.\ Models Methods Appl.\ Sci.\ \textbf{31} (2021), 2345--2390.

\bibitem{ShmaloUDC}
Y.~Shmalo,
\emph{Stability of accuracy for the training of DNNs via the uniform doubling
condition}, Ann.\ Math.\ Artif.\ Intell.\ (2024),
DOI 10.1007/s10472-023-09919-1.

\bibitem{companion2}
Y.~Shmalo, \emph{The uniform doubling condition holds with high probability for
independent samples}, manuscript, 2026.

\bibitem{companion3}
Y.~Shmalo, \emph{The critical scaling of the uniform doubling condition: smooth
bodies versus truncations}, manuscript, 2026.

\bibitem{companion4}
Y.~Shmalo, \emph{The logarithm is necessary: exact critical scaling of the
uniform doubling condition on the ball}, manuscript, 2026.

\bibitem{companion11}
Y.~Shmalo, \emph{The interpolation law for the critical scaling of the uniform
doubling condition}, manuscript, 2026.

\bibitem{companion6}
Y.~Shmalo, \emph{The edge exponent of the uniform doubling condition and an
interpolation conjecture}, manuscript, 2026.

\bibitem{companion9}
Y.~Shmalo, \emph{The interpolation law for the uniform doubling condition in
general position}, manuscript, 2026.

\end{thebibliography}

\end{document}
